\documentclass[11pt,a4paper]{article}
\usepackage[utf8]{inputenc}
\usepackage[vietnamese]{babel}
\usepackage{amsmath,amssymb}
\usepackage{geometry}
\geometry{a4paper, total={180mm,260mm}, left=15mm, right=15mm, top=18mm, bottom=18mm}
\usepackage{tikz}
\usetikzlibrary{shapes.geometric, arrows.meta, positioning, calc}
\usepackage{tcolorbox}
\tcbuselibrary{skins,breakable}
\usepackage{xcolor}
\usepackage{multicol}
\usepackage{array}
\usepackage{fontawesome}
\usepackage{fancyhdr}

\pagestyle{fancy}
\fancyhf{}
\renewcommand{\headrulewidth}{0.5pt}
\renewcommand{\footrulewidth}{0.5pt}
\fancyhead[L]{\small\bfseries\color{teal}\faBook\ SỔ TAY KIẾN THỨC TOÁN 3}
\fancyhead[R]{\small\bfseries\color{teal}Lớp Toán Thầy Lĩnh -- loptoanthaylinh.vn}
\fancyfoot[C]{\small\color{gray}-- Sổ tay kiến thức: Bài 78. Ôn tập phép nhân, phép chia trong phạm vi 100 000 --}

\definecolor{primary}{HTML}{005F73}
\definecolor{secondary}{HTML}{0A9396}
\definecolor{accent}{HTML}{EE9B00}
\definecolor{danger}{HTML}{AE2012}
\definecolor{boxbg}{HTML}{F4F6F8}

\begin{document}

% --- BANNER TIÊU ĐỀ ---
\begin{tcolorbox}[
  enhanced,
  colback=white,
  colframe=primary,
  arc=3mm,
  boxrule=1.2pt,
  left=4mm, right=4mm, top=2.5mm, bottom=2.5mm,
  drop shadow=black!10!white
]
  \centering
  {\Large\bfseries\color{primary}\faBookmark\ BÀI 78: ÔN TẬP NHÂN CHIA TRONG PHẠM VI 100 000}\\[1.5mm]
  {\small\itshape\color{black!75}Ôn tập cuối năm: Nhân chia số có đến 5 chữ số, xử lý chữ số 0 ở thương, bài toán rút về đơn vị trên 1 trang A4}
\end{tcolorbox}

\vspace{-2mm}


\noindent
\begin{minipage}[t]{0.49\textwidth}\vspace{0pt}
\begin{tcolorbox}[
  enhanced,
  colback=blue!4!white,
  colframe=blue!75!black,
  arc=2mm, boxrule=0.8pt,
  left=2mm, right=2mm, top=1.8mm, bottom=1.8mm,
  title={\footnotesize\bfseries\color{white}\faCalculator\ 1. Quy tắc Nhân và Chia cốt lõi}
]
{\footnotesize
\begin{itemize}\setlength{\itemsep}{1.5pt}\setlength{\leftskip}{-2.5mm}
  \item \textbf{Phép nhân:} Nhân lần lượt từ phải sang trái, nhân trước cộng nhớ sau.
  \item \textbf{Phép chia:} Chia lần lượt từ trái sang phải qua 3 nhịp: Chia $\to$ Nhân $\to$ Trừ.
  \item Khi hạ một chữ số mà số bị chia nhỏ hơn số chia $\implies$ phải ghi chữ số $\mathbf{0}$ vào thương!
  \item Số dư luôn nhỏ hơn số chia ($r < d$).
\end{itemize}
}
\end{tcolorbox}
\end{minipage}
\hfill
\begin{minipage}[t]{0.49\textwidth}\vspace{0pt}
\begin{tcolorbox}[
  enhanced,
  colback=teal!4!white,
  colframe=teal!75!black,
  arc=2mm, boxrule=0.8pt,
  left=2mm, right=2mm, top=1.8mm, bottom=1.8mm,
  title={\footnotesize\bfseries\color{white}\faSliders\ 2. Bài toán Rút về đơn vị và Thứ tự tính}
]
{\footnotesize
\begin{itemize}\setlength{\itemsep}{1.5pt}\setlength{\leftskip}{-2.5mm}
  \item \textbf{Bài toán Rút về đơn vị:}\\
  Bước 1: Tìm giá trị 1 phần (Phép chia $: n$).\\
  Bước 2: Tìm giá trị nhiều phần (Phép nhân $\times k$).
  \item \textbf{Thứ tự thực hiện biểu thức:}
  $$\text{Trong ngoặc trước} \to \text{Nhân chia trước} \to \text{Cộng trừ sau}$$
  Cùng cấp làm từ trái sang phải.
\end{itemize}
}
\end{tcolorbox}
\end{minipage}

\vspace{-2.5mm}
% --- KHỐI 3: SƠ ĐỒ TƯ DUY TIKZ ---
\noindent
\begin{tcolorbox}[
  enhanced,
  colback=orange!4!white,
  colframe=orange!80!black,
  arc=2mm, boxrule=0.8pt,
  left=2mm, right=2mm, top=1.5mm, bottom=1.5mm,
  title={\footnotesize\bfseries\color{white}\faSitemap\ SƠ ĐỒ TƯ DUY: HỆ THỐNG NHÂN CHIA ĐẠI TỔNG KẾT TOÁN 3}
]
\centering
\begin{tikzpicture}[
  box/.style={draw, rounded corners=2pt, align=center, font=\scriptsize\bfseries, inner sep=2.5pt, minimum height=0.75cm},
  arr/.style={-{Stealth[length=2mm]}, thick}
]
  \node[box, fill=purple!15, draw=purple!80!black, text width=2.6cm] (d1) at (0, 0) {BẢNG CỬU CHƯƠNG\\Bảng 2 đến 9\\Thuộc lòng tuyệt đối};
  \node[box, fill=teal!15, draw=teal!80!black, text width=2.6cm] (d2) at (3.5, 0) {ĐẶT TÍNH NHÂN CHIA\\Số có $2, 3, 4, 5$ chữ số\\với số có 1 chữ số};
  \node[box, fill=blue!12, draw=blue!80!black, text width=2.6cm] (d3) at (7.0, 0) {RÚT VỀ ĐƠN VỊ\\Bước 1: Tìm 1 phần ($:$)\\Bước 2: Tìm $k$ phần ($\times$)};
  \node[box, fill=orange!15, draw=orange!80!black, text width=2.6cm] (d4) at (10.5, 0) {BIỂU THỨC SỐ\\Quy tắc ưu tiên\\Tính nhanh phân phối};

  \draw[arr, purple!80!black] (d1) -- (d2);
  \draw[arr, teal!80!black] (d2) -- (d3);
  \draw[arr, blue!80!black] (d3) -- (d4);
\end{tikzpicture}
\end{tcolorbox}

\vspace{-2.5mm}
% --- KHỐI 4: VÍ DỤ & MẸO TRÁNH BẪY ---
\noindent
\begin{minipage}[t]{0.49\textwidth}\vspace{0pt}
\begin{tcolorbox}[
  enhanced,
  colback=green!4!white,
  colframe=green!60!black,
  arc=2mm, boxrule=0.8pt,
  left=2mm, right=2mm, top=1.8mm, bottom=1.8mm,
  title={\footnotesize\bfseries\color{white}\faCheckCircle\ Ví dụ thực hành chuẩn mực}
]
{\footnotesize
\textbf{1. Đặt tính chia:}\\
$42036 : 6 = \mathbf{7006}$.\\
-- $42 : 6 = 7$; hạ $0$, $0 : 6 = \mathbf{0}$; hạ $3$, $3 < 6 \implies \mathbf{0}$; hạ $6$ được $36$, $36 : 6 = 6$.\\
\textbf{2. Tính nhanh giá trị biểu thức:}\\
$125 \times 6 + 125 \times 4 = 125 \times (6 + 4) = 125 \times 10 = \mathbf{1250}$.
}
\end{tcolorbox}
\end{minipage}
\hfill
\begin{minipage}[t]{0.49\textwidth}\vspace{0pt}
\begin{tcolorbox}[
  enhanced,
  colback=red!4!white,
  colframe=red!70!black,
  arc=2mm, boxrule=0.8pt,
  left=2mm, right=2mm, top=1.8mm, bottom=1.8mm,
  title={\footnotesize\bfseries\color{white}\faExclamationTriangle\ Mẹo tránh bẫy trắc nghiệm}
]
{\footnotesize
\begin{itemize}\setlength{\itemsep}{1.5pt}\setlength{\leftskip}{-2.5mm}
  \item \textbf{Bẫy quên cả hai chữ số 0 ở thương:}\\
  Ở phép chia $42\,036 : 6$, nếu quên số 0 thì ra $76$! Thử lại: $7006 \times 6 = 42\,036$ mới là đáp số chuẩn xác.
  \item Cảnh giác với các phép chia có dư: số dư luôn phải nhỏ hơn số chia!
\end{itemize}
}
\end{tcolorbox}
\end{minipage}


\end{document}
